% Introduction to the queen vs queen chapter, for the 1500-1599 book only.
% There is no shared texts/queen-vs-queen.tex; the theory is in texts/queen-vs-queen/one-pawn-vs-none.tex.

\section*{What to look for}
Queen endings often involve long variations, and the clearly right move is hard to find. The main
difficulty is the perpetual check:\index{Perpetual check} almost every move gives or meets a check,
and you often cannot tell whether the checks will ever stop. That makes these endings hard to
assess. The few rules we know for queen and pawn against queen are given in
\xref{Section}{queen-and-pawn}.

Eliminate the wrong moves. Since the right move is rarely obvious, the
\term{process-of-elimination}\index{Practical play!Process of elimination} is the main tool: rule out
the moves that allow a perpetual check or lose a pawn, and choose among the rest
(\xref{Diagram}{saim-aslan-92w} and \xref{Diagram}{sabov-georgiev-70w}). From a practical point of
view, it is sometimes enough to keep the position going and hope that the opponent goes wrong.

Check along the diagonals. Checks along the diagonals reach a perpetual check more often than
checks along files and ranks,\index{Queen endings!Check along the diagonal} and the defender's king is
safest where the attacker cannot answer a check with a check (\xref{Diagram}{sabov-georgiev-71w}
and \xref{Diagram}{duman-yildirak-80w}).

The defender has good chances. Thanks to perpetual check, the defender can often hold even
two or three pawns down. When the king has no moves left, look for a desperado\index{Desperado}
(\xref{Diagram}{silvestre-dias-68b} and \xref{Diagram}{hamoonga-salimmoussa-52b}).

Think twice before entering a queen ending. For the same reason, going into a queen ending is
unwise if you have a better alternative. A pawn ending, where the result can be counted, is often
the safer way to win\index{Transposing into a pawn ending} (\xref{Diagram}{bolu-koc-53b} and
\xref{Diagram}{gabatshwarwerefilwe-mongeli-76b}).
